\documentclass[10pt,a4paper]{amsart}
\usepackage[margin=19mm]{geometry}
\usepackage{amsmath,amssymb,mathtools}
\usepackage[scr=rsfs,cal=euler]{mathalfa}
\usepackage{xfrac}
\usepackage[unicode,hidelinks,bookmarks=false]{hyperref}
\newcommand{\ie}{i.e.\ }
\newcommand{\cf}{cf.\ }
\newcommand{\resp}{respectively\ }
\newcommand{\Subsection}{\subsection}
\providecommand{\nobreakdash}{\nobreak\hbox{-}}
\setlength{\emergencystretch}{2em}
\multlinegap0pt
\usepackage[shortlabels]{enumitem}
\usepackage{amssymb}
\usepackage{mathtools}
\usepackage{bm}
\usepackage{xcolor}
\newtheorem{problem}{Problem}
\newtheorem{theorem}{Theorem}
\newtheorem{remark}{Remark}
\newcommand{\R}{\mathbb{R}}
\newcommand{\dd}{\,\mathrm{d}}
\title{\bfseries Age-Structured Harvesting Models: A Structural Comparison of Rate-Control and Effort-Control Optimality Systems}
\author{Jiguang Yu, Louis Shuo Wang, Ye Liang}
\date{Source: arXiv:2603.09767v3 (CC BY 4.0)}
\begin{document}
\maketitle

\noindent\emph{Source and licence.} \bfseries Age-Structured Harvesting Models: A Structural Comparison of Rate-Control and Effort-Control Optimality Systems. Version arXiv:2603.09767v3 (CC BY 4.0), distributed by arXiv under the Creative Commons Attribution 4.0 International licence, http://creativecommons.org/licenses/by/4.0/, as recorded in the arXiv licence metadata for this exact version and shown on the abstract page https://arxiv.org/abs/2603.09767v3. Full licence notice: \texttt{licenses/arxiv-2603.09767-CC-BY-4.0.txt}.\par\smallskip

\noindent\textit{Editorial scope.} Counted unit: the article's theorem thm:necessary\_PR (source0.tex lines 769--815). The article's complete original proof of it is delivered with it (source0.tex lines 817--854, one \texttt{proof} environment of 38 lines). No mathematics is written, repaired or re-derived: the statement and the proof are the article's own lines, in the article's own notation, and no retained line is substituted or re-flowed.  Retained uncounted as the source-local context the proof needs: lines 508--524; lines 545--554; lines 650--670; lines 737--749; lines 761--764; lines 1204--1207.  Nothing was omitted from any retained span, so the package carries no omission record.  No retained line contains a citation, so the wrapper carries no bibliography and no bibliography entry.  No retained line contains a figure environment, an image-inclusion command or drawing code, so no image is delivered and the package declares no asset dependency.  Editorial formatting only, in addition: an AMS \texttt{amsart} preamble that restates the article's own declarations and macros used by the retained lines, namely the environments \texttt{problem}, \texttt{theorem}, \texttt{remark} on separate counters, together with the standard packages those lines need.  The retained environments print in this document as Problem 1, Theorem 1, Remark 1 in order of appearance, whereas the article prints the same items with the numbering of its own full document.  The complete native source of the article as distributed in the arXiv e-print for this exact version is bundled unchanged as \texttt{sources/196085/source0.tex}.
\begin{problem}\label{prob:PR}
Maximise
\begin{equation}\label{eq:PR_obj}
J(u,p)=\int_0^\infty e^{-rt}
\left(\int_0^A c(t,a)\,u(t,a)\dd a-k(t)\,p(t)\right)\dd t
\end{equation}
over all state-admissible control pairs $(u,p)\in\mathcal{U}_{R,\mathrm{ad}}$,
where for each such pair the state $x=S(u,p)$ is the unique mild solution of
\begin{align}
\partial_t x+\partial_a x&=-\mu(t,a)x-u(t,a), && t>0,\ a\in(0,A),
\label{eq:PR_state}\\
x(t,0)&=p(t), && t\ge0, \label{eq:PR_bc}\\
x(0,a)&=x_0(a), && a\in[0,A]. \label{eq:PR_ic}
\end{align}
Here $r>0$ is the discount rate, $c(t,a)\ge0$ is the unit harvesting value,
and $k(t)\ge0$ is the boundary inflow cost.
\end{problem}

\begin{theorem}[Existence for the rate-control problem]
\label{thm:existence_PR}
Assume \textup{(A1)--(A3)}, $r>0$, $c\in L^\infty((0,\infty)\times(0,A))$,
$k\in L^\infty(0,\infty)$, and that the control-to-state map
$S:\mathcal{U}_R\to L^1_{\mathrm{loc}}([0,\infty);L^1(0,A))$ is affine and
sequentially weak-$\ast$-to-$L^1_{\mathrm{loc}}$ continuous. If
$\mathcal{U}_{R,\mathrm{ad}}\neq\varnothing$, then Problem~\ref{prob:PR}
admits at least one optimal control
$(u^\ast,p^\ast)\in\mathcal{U}_{R,\mathrm{ad}}$.
\end{theorem}

\begin{remark}[Measure-valued multipliers]
\label{rem:measure_multiplier}
For pointwise pure-state constraints such as $x(t,a)\ge0$, the natural
Lagrange multiplier is in general a nonnegative Radon measure
$\eta\in\mathcal{M}_+([0,\infty)\times[0,A])$ supported on the active set
$\{x^\ast=0\}$, rather than an $L^1_{\mathrm{loc}}$ density; this is the
classical phenomenon for state-constrained control. Assuming
$\eta\in L^1_{\mathrm{loc}}$ in (R5) is therefore a \emph{simplifying
regularity hypothesis}, valid when the active set is ``fat'' enough that the
multiplier admits a density, and convenient because it keeps all terms in
the integration-by-parts formula classical. The entire derivation below goes
through verbatim with $\eta$ a measure provided the products
$e^{-rt}\lambda\dd\eta$ and the pairing $\langle\eta,\delta x\rangle$ are
interpreted in the measure sense and $\delta x$ is taken continuous in
$(t,a)$; the adjoint equation \eqref{eq:adjoint_theorem_PR} then holds in the
sense of measures, with $\eta$ contributing a singular source on the active
set. We record the $L^1$ form for readability and flag the measure-valued
generalisation as the rigorous formulation. This point is especially
relevant for Model~R, where additive removal can activate the constraint on
specific age classes.
\end{remark}

\begin{align}
\delta\mathcal L
&=\int_0^\infty\!\!\int_0^A e^{-rt}\bigl(c-\lambda\bigr)\delta u\dd a\dd t
+\int_0^\infty e^{-rt}\bigl(\lambda(t,0)-k\bigr)\delta p\dd t
\notag\\
&\quad+\int_0^\infty\!\!\int_0^A e^{-rt}
\Bigl(\partial_t\lambda+\partial_a\lambda-(r+\mu)\lambda+\eta\Bigr)
\delta x\dd a\dd t
\notag\\
&\quad-\int_0^\infty e^{-rt}\lambda(t,A)\delta x(t,A)\dd t
-\lim_{T\to\infty}\int_0^A e^{-rT}\lambda(T,a)\delta x(T,a)\dd a.
\label{eq:pf_deltaL_3}
\end{align}

\begin{equation}\label{eq:pf_deltaL_reduced}
\delta\mathcal L=\int_0^\infty\!\!\int_0^A e^{-rt}(c-\lambda)\delta u\dd a\dd t
+\int_0^\infty e^{-rt}(\lambda(t,0)-k)\delta p\dd t,
\end{equation}

\begin{equation}\label{eq:auto_adjoint_formula}
\lambda(a)=\int_a^A\exp\!\left(-\int_a^s\bigl(r+\mu(\xi)\bigr)\dd\xi\right)
\eta(s)\dd s,\qquad 0\le a\le A.
\end{equation}

\begin{theorem}[Conditional necessary optimality system]
\label{thm:necessary_PR}
Let $(u^\ast,p^\ast,x^\ast)$ be an optimal solution of
Problem~\ref{prob:PR}, and assume \textup{(R1)--(R7)}. Then there exist a
current-value adjoint $\lambda:[0,\infty)\times[0,A]\to\R$ and a current-value
multiplier $\eta:[0,\infty)\times[0,A]\to[0,\infty)$ (interpreted as a
measure in the sense of Remark~\ref{rem:measure_multiplier} if (R5) is
relaxed) such that:
\begin{enumerate}[leftmargin=2.6em]
    \item[\textup{(i)}] \textbf{Adjoint equation.}
    \begin{equation}\label{eq:adjoint_theorem_PR}
    \partial_t\lambda+\partial_a\lambda=(r+\mu)\lambda-\eta,
    \qquad t>0,\ a\in(0,A).
    \end{equation}
    \item[\textup{(ii)}] \textbf{Terminal-age and transversality conditions.}
    \begin{equation}\label{eq:adjoint_terminal_theorem_PR}
    \lambda(t,A)=0\ (t\ge0),\qquad
    \lim_{T\to\infty}\int_0^A e^{-rT}\lambda(T,a)\delta x(T,a)\dd a=0
    \end{equation}
    for every admissible first variation $\delta x$.
    \item[\textup{(iii)}] \textbf{Switching for the distributed control $u$.}
    With $\frac{\partial I}{\partial u}(t,a):=e^{-rt}(c(t,a)-\lambda(t,a))$,
    \begin{equation}\label{eq:u_switch}
    \frac{\partial I}{\partial u}\le0\ \text{where }u^\ast=0,\quad
    \frac{\partial I}{\partial u}=0\ \text{where }0<u^\ast<u_{\max},\quad
    \frac{\partial I}{\partial u}\ge0\ \text{where }u^\ast=u_{\max}.
    \end{equation}
    \item[\textup{(iv)}] \textbf{Switching for the boundary control $p$.}
    With $\frac{\partial I}{\partial p}(t):=e^{-rt}(\lambda(t,0)-k(t))$,
    \begin{equation}\label{eq:p_switch}
    \frac{\partial I}{\partial p}\le0\ \text{where }p^\ast=0,\quad
    \frac{\partial I}{\partial p}=0\ \text{where }0<p^\ast<p_{\max},\quad
    \frac{\partial I}{\partial p}\ge0\ \text{where }p^\ast=p_{\max}.
    \end{equation}
    \item[\textup{(v)}] \textbf{Complementary slackness.}
    $\eta\ge0$, $x^\ast\ge0$, and $\eta\,x^\ast=0$ a.e.\ (equivalently,
    $\operatorname{supp}\eta\subseteq\{x^\ast=0\}$ in the measure-valued
    formulation).
\end{enumerate}
Moreover, given the multiplier $\eta$ and the data, the adjoint $\lambda$ is
\emph{unique}: it is the unique solution of the backward transport problem
\eqref{eq:adjoint_theorem_PR} with terminal condition
$\lambda(t,A)=0$, integrated along characteristics from $a=A$ downward
(cf.\ the explicit stationary formula \eqref{eq:auto_adjoint_formula}). The
multiplier $\eta$ is in turn determined on the active set by complementary
slackness.
\end{theorem}

\begin{proof}
Existence of an optimal pair is Theorem~\ref{thm:existence_PR}. By (R3),
(R5), (R6) the first variation of the current-value Lagrangian is
\eqref{eq:pf_deltaL_3}, and choosing $\lambda$, the terminal-age condition,
and the transversality condition as above reduces it to
\eqref{eq:pf_deltaL_reduced}. This proves (i), (ii).

For (iii)--(iv) we use the normal-cone characterisation of the box
constraints. At the optimum,
$\delta\mathcal L\le0$ for every feasible direction $(\delta u,\delta p)$,
i.e.\ for every $(\delta u,\delta p)$ in the tangent cone to the box
$[0,u_{\max}]\times[0,p_{\max}]$ at $(u^\ast,p^\ast)$. Equivalently, the
gradient pair $(\frac{\partial I}{\partial u},\frac{\partial I}{\partial p})$
lies in the normal cone $N_{\mathcal U_R}(u^\ast,p^\ast)$. For the pointwise
box constraint $0\le u^\ast(t,a)\le u_{\max}$ this normal cone decomposes
pointwise as
\[
N(t,a)=\begin{cases}
(-\infty,0], & u^\ast(t,a)=0,\\
\{0\}, & 0<u^\ast(t,a)<u_{\max},\\
[0,\infty), & u^\ast(t,a)=u_{\max},
\end{cases}
\]
which is exactly \eqref{eq:u_switch}; the same argument applied to
$p^\ast$ gives \eqref{eq:p_switch}. Concretely: where $0<u^\ast<u_{\max}$ both
signs of $\delta u$ are admissible, forcing
$\frac{\partial I}{\partial u}=0$; where $u^\ast=0$ only $\delta u\ge0$ is
admissible, forcing $\frac{\partial I}{\partial u}\le0$; where
$u^\ast=u_{\max}$ only $\delta u\le0$ is admissible, forcing
$\frac{\partial I}{\partial u}\ge0$.

Part (v) is the constraint qualification (R7): the pointwise state constraint
$x^\ast\ge0$ is enforced by a multiplier $\eta\ge0$ satisfying complementary
slackness $\eta\,x^\ast=0$. Uniqueness of $\lambda$ follows because
\eqref{eq:adjoint_theorem_PR} with $\lambda(t,A)=0$ is a backward transport
equation with a one-sided (terminal-age) condition, integrable uniquely along
characteristics.
\end{proof}

\end{document}
